\chapter{Lattice vibrations and electron--phonon coupling}
\setwlchapter{phonons}
\section{Quantizing a displacement}
For a monatomic chain with spring constant $K$ and mass $M$,
$\omega_q=2\sqrt{K/M}\,|\sin(qa/2)|$. For $\omega_q>0$ the normal
coordinate is proportional to $(b_q+b^\dagger_{-q})/\sqrt{2M\omega_q}$;
the adjoint momentum involves their difference. Pairing $q$ with $-q$
keeps the real-space displacement Hermitian. The centre-of-mass
translation at $q=0$ is a free coordinate, not an oscillator for which
one may substitute $\omega=0$ into this normalization.

\section{Density and bond couplings}
A local Holstein term $g(n_i-n_0)(b_i+b_i^\dagger)$ describes a change
of site energy. On a bond, modulation of the hopping instead gives
$\alpha(u_{i+1}-u_i)\sum_\sigma(c^\dagger_{i\sigma}c_{i+1,\sigma}+\hc)$.
Uniform displacement drops out of the bond coupling. Fourier transforming
it supplies a factor $e^{iqa}-1$ that a local optical-mode vertex does
not have. In either case a phonon carries the momentum transferred
between electron states.
\casestudy{holstein}{Checking the momentum of the emitted phonon}

\section{Displaced oscillator and induced attraction}
For a fixed electronic occupation put $B=b+g(n-n_0)/\omega_0$. Then
\[
 \omega_0 b^\dagger b+g(n-n_0)(b+b^\dagger)
 =\omega_0 B^\dagger B-g^2(n-n_0)^2/\omega_0 .
\]
The ground energy depends on charge even when the bare phonon has no
dispersion. The second difference between energies of charge 2, 1, 0
receives a shift $-2g^2/\omega_0$. In a hopping system a change of
occupation also changes the displacement; the transformed hopping
carries displacement operators. Replacing them by their phonon-vacuum
expectations is an approximation. In an atomic charge sector the
completion of the square is exact.
\casestudy{lang-firsov}{Convergence with the number of phonon quanta}
The displaced vacuum contains a Poisson distribution in the original
oscillator basis with mean $[g(n-n_0)/\omega_0]^2$. The largest
occupied number in a truncated eigenvector is a useful convergence
diagnostic in addition to its energy. On a dimer check the convergence
of the bonding--antibonding splitting itself: it can be much smaller
than the shift in an individual energy.

\subsection{How a phonon cloud changes the hopping}
Take one spinless electron and two local Einstein oscillators on a
symmetric dimer. In the absence of tunnelling the two localized
ground states are degenerate. Moving the electron interchanges which
oscillator is displaced. Their phonon wave functions have overlap
$e^{-(g/\omega_0)^2}$: each of two oscillators contributes
$e^{-(g/\omega_0)^2/2}$. Degenerate perturbation theory gives a
splitting $2t e^{-(g/\omega_0)^2}$ to first order in $t$, with coupling
$g n_j(b_j+b_j^\dagger)$ and real positive hopping. This splitting is
a different observable from the single-site relaxation energy.

Build the matrix directly as a check on the phonon convention. The
tensor order is electron site, phonon on the left, phonon on the
right; $t$ and $g$ are in units of $\omega_0$.
\begin{wl}
split[n_,t_,g_] := Module[{id,nop,x,ph,h},
  id = IdentityMatrix[n+1]; nop = DiagonalMatrix[Range[0,n]];
  x = Normal[SparseArray[{Band[{1,2}] -> Sqrt[Range[n]],
                  Band[{2,1}] -> Sqrt[Range[n]]},{n+1,n+1}]];
  ph = KroneckerProduct[nop,id] + KroneckerProduct[id,nop];
  h = KroneckerProduct[IdentityMatrix[2],ph] -
      t KroneckerProduct[PauliX,id,id] +
      g (KroneckerProduct[DiagonalMatrix[{1,0}],x,id] +
         KroneckerProduct[DiagonalMatrix[{0,1}],id,x]);
  Differences[Take[Sort[Eigenvalues[N[h]]],2]][[1]]];
TableForm[Table[{n,split[n,0.01,0.6]/0.02, Exp[-0.6^2]}, {n,{2,4,6}}],
 TableHeadings -> {None,{"Cutoff","Splitting / 2t","Overlap"}}]
\end{wl}
Converging the cutoff at fixed $t$ approaches the finite-hopping
answer. Reducing $t$ at fixed converged cutoff is the separate test
of first-order perturbation theory.

\section{Virtual phonons and coupled modes}
At imaginary transferred frequency $\nu$, eliminating a harmonic
phonon coupled linearly to density yields
$V(q,i\nu)=-2|g_q|^2\omega_q/(\nu^2+\omega_q^2)$ for the dimensionless
coordinate $b+b^\dagger$. Only for $|\nu|\ll\omega_q$ can this be
replaced by a static attraction. A canonical transformation in powers
of $g$ must also avoid resonances between real electronic transitions
and a phonon; there the eliminated state belongs in the active space.
\casestudy{squeezing}{Stability of two coupled harmonic modes}

For an orbital doublet coupled to two degenerate distortions the
linear Jahn--Teller Hamiltonian at fixed coordinates contains
$g(Q_1\tau_z+Q_2\tau_x)$. Its lower adiabatic surface is
$K(Q_1^2+Q_2^2)/2-|g|\sqrt{Q_1^2+Q_2^2}$. The continuum of classical
minima does not imply a static symmetry-breaking expectation value in
a finite quantum molecule: the kinetic energy of the nuclei mixes
distortions. A finite two-mode Fock calculation can retain both.
\begin{exercise}[Relaxation energy]\label{ex:jt}
Find the Jahn--Teller minimum radius and lowering of the energy.
Relate the latter to the displaced-oscillator result for a single
distortion quantized with $Q=(b+b^\dagger)/\sqrt{2M\omega}$.
\end{exercise}
